\documentclass[11pt,preprint]{elsarticle}

\usepackage{lmodern}
%%%% My spacing
\usepackage{setspace}
\setstretch{1.2}
\DeclareMathSizes{12}{14}{10}{10}

% Wrap around which gives all figures included the [H] command, or places it "here". This can be tedious to code in Rmarkdown.
\usepackage{float}
\let\origfigure\figure
\let\endorigfigure\endfigure
\renewenvironment{figure}[1][2] {
    \expandafter\origfigure\expandafter[H]
} {
    \endorigfigure
}

\let\origtable\table
\let\endorigtable\endtable
\renewenvironment{table}[1][2] {
    \expandafter\origtable\expandafter[H]
} {
    \endorigtable
}


\usepackage{ifxetex,ifluatex}
\usepackage{fixltx2e} % provides \textsubscript
\ifnum 0\ifxetex 1\fi\ifluatex 1\fi=0 % if pdftex
  \usepackage[T1]{fontenc}
  \usepackage[utf8]{inputenc}
\else % if luatex or xelatex
  \ifxetex
    \usepackage{mathspec}
    \usepackage{xltxtra,xunicode}
  \else
    \usepackage{fontspec}
  \fi
  \defaultfontfeatures{Mapping=tex-text,Scale=MatchLowercase}
  \newcommand{\euro}{€}
\fi

\usepackage{amssymb, amsmath, amsthm, amsfonts}

\def\bibsection{\section*{References}} %%% Make "References" appear before bibliography


\usepackage[numbers]{natbib}

\usepackage{longtable}
\usepackage[margin=2.3cm,bottom=2cm,top=2.5cm, includefoot]{geometry}
\usepackage{fancyhdr}
\usepackage[bottom, hang, flushmargin]{footmisc}
\usepackage{graphicx}
\numberwithin{equation}{section}
\numberwithin{figure}{section}
\numberwithin{table}{section}
\setlength{\parindent}{0cm}
\setlength{\parskip}{1.3ex plus 0.5ex minus 0.3ex}
\usepackage{textcomp}
\renewcommand{\headrulewidth}{0.2pt}
\renewcommand{\footrulewidth}{0.3pt}

\usepackage{array}
\newcolumntype{x}[1]{>{\centering\arraybackslash\hspace{0pt}}p{#1}}

%%%%  Remove the "preprint submitted to" part. Don't worry about this either, it just looks better without it:

 \def\tightlist{} % This allows for subbullets!

\usepackage{hyperref}
\hypersetup{breaklinks=true,
            bookmarks=true,
            colorlinks=true,
            citecolor=blue,
            urlcolor=blue,
            linkcolor=blue,
            pdfborder={0 0 0}}


% The following packages allow huxtable to work:
\usepackage{siunitx}
\usepackage{multirow}
\usepackage{hhline}
\usepackage{calc}
\usepackage{tabularx}
\usepackage{booktabs}
\usepackage{caption}


\newenvironment{columns}[1][]{}{}

\newenvironment{column}[1]{\begin{minipage}{#1}\ignorespaces}{%
\end{minipage}
\ifhmode\unskip\fi
\aftergroup\useignorespacesandallpars}

\def\useignorespacesandallpars#1\ignorespaces\fi{%
#1\fi\ignorespacesandallpars}

\makeatletter
\def\ignorespacesandallpars{%
  \@ifnextchar\par
    {\expandafter\ignorespacesandallpars\@gobble}%
    {}%
}
\makeatother


% definitions for citeproc citations
\NewDocumentCommand\citeproctext{}{}
\NewDocumentCommand\citeproc{mm}{%
\href{\#cite.\detokenize{#1}}{#2}\nocite{#1}}

\makeatletter
% allow citations to break across lines
\let\@cite@ofmt\@firstofone
% avoid brackets around text for \cite:
\def\@biblabel#1{}
\def\@cite#1#2{{#1\if@tempswa , #2\fi}}
\makeatother
\newlength{\cslhangindent}
\setlength{\cslhangindent}{1.5em}
\newlength{\csllabelwidth}
\setlength{\csllabelwidth}{3em}
\newenvironment{CSLReferences}[2] % #1 hanging-indent, #2 entry-spacing
{\begin{list}{}{%
	\setlength{\itemindent}{0pt}
	\setlength{\leftmargin}{0pt}
	\setlength{\parsep}{0pt}
	% turn on hanging indent if param 1 is 1
	\ifodd #1
	\setlength{\leftmargin}{\cslhangindent}
	\setlength{\itemindent}{-1\cslhangindent}
	\fi
	% set entry spacing
	\setlength{\itemsep}{#2\baselineskip}}}
{\end{list}}

\usepackage{calc}
\newcommand{\CSLBlock}[1]{\hfill\break\parbox[t]{\linewidth}{\strut\ignorespaces#1\strut}}
\newcommand{\CSLLeftMargin}[1]{\parbox[t]{\csllabelwidth}{\strut#1\strut}}
\newcommand{\CSLRightInline}[1]{\parbox[t]{\linewidth - \csllabelwidth}{\strut#1\strut}}
\newcommand{\CSLIndent}[1]{\hspace{\cslhangindent}#1}


\urlstyle{same}  % don't use monospace font for urls
\setlength{\parindent}{0pt}
\setlength{\parskip}{6pt plus 2pt minus 1pt}
\setlength{\emergencystretch}{3em}  % prevent overfull lines
\setcounter{secnumdepth}{5}

%%% Use protect on footnotes to avoid problems with footnotes in titles
\let\rmarkdownfootnote\footnote%
\def\footnote{\protect\rmarkdownfootnote}
\IfFileExists{upquote.sty}{\usepackage{upquote}}{}

%%% Include extra packages specified by user

%%% Hard setting column skips for reports - this ensures greater consistency and control over the length settings in the document.
%% page layout
%% paragraphs
\setlength{\baselineskip}{12pt plus 0pt minus 0pt}
\setlength{\parskip}{12pt plus 0pt minus 0pt}
\setlength{\parindent}{0pt plus 0pt minus 0pt}
%% floats
\setlength{\floatsep}{12pt plus 0 pt minus 0pt}
\setlength{\textfloatsep}{20pt plus 0pt minus 0pt}
\setlength{\intextsep}{14pt plus 0pt minus 0pt}
\setlength{\dbltextfloatsep}{20pt plus 0pt minus 0pt}
\setlength{\dblfloatsep}{14pt plus 0pt minus 0pt}
%% maths
\setlength{\abovedisplayskip}{12pt plus 0pt minus 0pt}
\setlength{\belowdisplayskip}{12pt plus 0pt minus 0pt}
%% lists
\setlength{\topsep}{10pt plus 0pt minus 0pt}
\setlength{\partopsep}{3pt plus 0pt minus 0pt}
\setlength{\itemsep}{5pt plus 0pt minus 0pt}
\setlength{\labelsep}{8mm plus 0mm minus 0mm}
\setlength{\parsep}{\the\parskip}
\setlength{\listparindent}{\the\parindent}
%% verbatim
\setlength{\fboxsep}{5pt plus 0pt minus 0pt}



\begin{document}



\begin{frontmatter}  %

\title{Coffe Hub}

% Set to FALSE if wanting to remove title (for submission)




\author[Add1]{25424971}
\ead{25424971@sun.ac.za}








\vspace{1cm}





\vspace{0.5cm}

\end{frontmatter}

\setcounter{footnote}{0}



%________________________
% Header and Footers
%%%%%%%%%%%%%%%%%%%%%%%%%%%%%%%%%
\pagestyle{fancy}
\chead{}
\rhead{}
\lfoot{}
\rfoot{\footnotesize Page \thepage}
\lhead{}
%\rfoot{\footnotesize Page \thepage } % "e.g. Page 2"
\cfoot{}

%\setlength\headheight{30pt}
%%%%%%%%%%%%%%%%%%%%%%%%%%%%%%%%%
%________________________

\headsep 35pt % So that header does not go over title




\section{Introduction}\label{introduction}

This section explores the potential coffee distributors for a coffee
shop in Stellenbosch. The cost (in USD), the roaster company, the roast
strength, and reviews of the coffees are considered.

\section{Average Rating and Price}\label{average-rating-and-price}

\subsection{Average Rating and Price by Roaster
Location}\label{average-rating-and-price-by-roaster-location}

The plot below indicates that Australia and England are clear outliers.
Although roasters in these countries have, on average, the highest
ratings, all the countries receive an average rating of 90 or above. I
would suggest that the prices disqualify the further consideration of
Australian and English roasters.

\begin{figure}[H]

{\centering \includegraphics{25424971Question1_files/figure-latex/unnamed-chunk-2-1} 

}

\caption{Average Rating and Price by Roaster Country.\label{Figure1}}\label{fig:unnamed-chunk-2}
\end{figure}

\subsection{Average Rating and Price by Roaster Location, excluding
Outliers}\label{average-rating-and-price-by-roaster-location-excluding-outliers}

Ratings cluster between 92 and 94, however, outliers in price remain.

\begin{figure}[H]

{\centering \includegraphics{25424971Question1_files/figure-latex/unnamed-chunk-3-1} 

}

\caption{Average Rating and Price by Location of Roaster without Outliers.\label{Figure1}}\label{fig:unnamed-chunk-3}
\end{figure}

\subsection{Average Rating and Price by Roast
Level}\label{average-rating-and-price-by-roast-level}

Light roast has the highest average rating, but also the highest average
cost. Dark has the lowest average rating and the lowest cost.
Medium-light and medium roast offer a middle-ground.

\begin{figure}[H]

{\centering \includegraphics{25424971Question1_files/figure-latex/unnamed-chunk-5-1} 

}

\caption{Average Rating and Price by Roast Level.\label{Figure1}}\label{fig:unnamed-chunk-5}
\end{figure}

\section{Indicator words}\label{indicator-words}

Wrangle the data using a function:

\subsection{Average Rating and Price by Review Indicator
Words}\label{average-rating-and-price-by-review-indicator-words}

This plot indicates that coffee brands that have reviews describing the
aroma and chocolatiness of the coffee are the best deal when it comes to
the ratings-price relationship. Subsequent analysis will be focused on
these observations.

\begin{figure}[H]

{\centering \includegraphics{25424971Question1_files/figure-latex/unnamed-chunk-7-1} 

}

\caption{Average Rating and Price by Review Indicator Words.\label{Figure1}}\label{fig:unnamed-chunk-7}
\end{figure}

\subsection{Boxplot by Country and Roast
Level}\label{boxplot-by-country-and-roast-level}

The goal is to trim the most expensive coffee products. Boxplot shows
the outliers in red. The boxplot also illustrates that Taiwan and the US
are the have the most roasters.

\begin{figure}[H]

{\centering \includegraphics{25424971Question1_files/figure-latex/unnamed-chunk-8-1} 

}

\caption{Price by Country and Roast level.\label{Figure1}}\label{fig:unnamed-chunk-8}
\end{figure}

\subsection{Price Distribution}\label{price-distribution}

The distribution in price aids in determining where a hard cap on the
price should be. At USD8 per 100 grams, there is a fall in the number of
coffee products sold at this price, and provides a reasonable cap.

\begin{figure}[H]

{\centering \includegraphics{25424971Question1_files/figure-latex/unnamed-chunk-9-1} 

}

\caption{Price Distribution.\label{Figure1}}\label{fig:unnamed-chunk-9}
\end{figure}

\subsection{Tile Plot}\label{tile-plot}

This plot depicts the average ratings by roaster and roast level once
the coffee products that have a cost larger than USD8 and roasters that
only supply one remaining product have been filtered out. The reasoning
is that a distributor can provide a diversified selction of products.

\begin{figure}[H]

{\centering \includegraphics{25424971Question1_files/figure-latex/unnamed-chunk-10-1} 

}

\caption{Average Rating by Roaster and Roast Level.\label{Figure1}}\label{fig:unnamed-chunk-10}
\end{figure}

\subsection{Prices and Ratings for
Roasters}\label{prices-and-ratings-for-roasters}

The plot below displays the price-ratings relationship for narrowd-down
roasters. Kakalove Cafe seems to be a promising supplier, with most of
its product prices falling the lowest, and receiving the highest some of
the highest ratings. They also produce a variety of products.

\begin{figure}[H]

{\centering \includegraphics{25424971Question1_files/figure-latex/unnamed-chunk-11-1} 

}

\caption{Ultimate Prices and Ratings.\label{Figure1}}\label{fig:unnamed-chunk-11}
\end{figure}

\bibliography{Tex/ref}





\end{document}
